%%% FIELD universal-recovery-proof
Fix an arbitrary member of \(\mathcal M_p\).  Thus \(p\geq 2\), every
\(\Sigma_c\) is positive definite, the weights are positive and sum to one,
the baseline factorization in \(\texttt{ass:baseline-gaussian-sem}\) holds,
and every intervention satisfies
\(\texttt{ass:perfect-single-node-surgeries}\) and
\(\texttt{ass:effective-interventions}\).  By
\(\texttt{lem:component-distinctness}\), \(P_{\mathrm{mix}}\) determines the
unordered collection of distinct covariance--weight pairs
\[
 \bigl\{(\Sigma_c,\pi_c):c\in C\bigr\}
\]
up to a simultaneous permutation of their component positions.  In
particular, inversion is legal and determines every
\(K_c=\Sigma_c^{-1}\).  We now construct the semantic explanation from this
observable collection and prove that the construction returns the generating
one.

First we recover the baseline position.  Every finite DAG has a root: if each
vertex had a parent, repeatedly choosing a parent would revisit a vertex and
produce a directed cycle.  Let \(r\) be a root.  Then
\(a_r=e_r\), and effectiveness gives \(\tau_r\neq\omega_r\).  Hence the
surgery identity gives
\[
 K_r-K_0
 =\left(\tau_r^{-1}-\omega_r^{-1}\right)e_re_r^{\top},
 \qquad \tau_r^{-1}-\omega_r^{-1}\neq0.                 \tag{1}
\]
Thus at least one coordinate-isolated pair is present in the observed
collection.  Conversely, by
\(\texttt{lem:coordinate-pair-characterization}\), every pair isolated at a
coordinate \(r\) consists exactly of the baseline and the intervention on the
root \(r\).  Since \(p\geq2\), the collection has \(p+1\geq3\) components, so
there is a third component.  The same lemma states that, of the two endpoints
of the isolated pair, exactly the baseline endpoint has its \(r\)-th marginal
variance equal to that of any third component.  All quantities in this test
are entries of observed covariance matrices.  It therefore orients every
coordinate-isolated pair toward one and the same position, namely the true
baseline position \(c_0\).  This proves existence and uniqueness of the
recovered \(c_0\).

For each remaining component form the observable set
\[
 S_c=\operatorname{rowsupp}(K_c-K_{c_0}),
 \qquad c\in C\setminus\{c_0\}.                         \tag{2}
\]
Temporarily use the true target label \(j\) for the component generated by
the intervention on \(j\).  The support conclusion of
\(\texttt{lem:surgery-identity}\), which also covers the allowed nonroot case
\(\tau_j=\omega_j\), gives
\[
 S_j=\{j\}\cup\operatorname{Pa}(j).                    \tag{3}
\]
Consequently the true component-to-target map is a bijection whose assigned
coordinate lies in the corresponding set in (2).

That support-compatible bijection is unique.  Indeed, let another such
bijection be given and, after identifying components by their true targets,
write it as a permutation \(f:V\to V\).  Equation (3) implies
\[
 f(j)\in\{j\}\cup\operatorname{Pa}(j)
 \quad\text{for every }j\in V.                          \tag{4}
\]
A finite acyclic graph admits a rank \(q:V\to\{1,\ldots,p\}\) that is
strictly increasing along each edge: recursively remove a root and assign
successive ranks.  Therefore (4) yields
\[
 q(f(j))\leq q(j),
\]
with strict inequality whenever \(f(j)\neq j\).  But bijectivity of \(f\)
also yields
\[
 \sum_{j\in V}q(f(j))=\sum_{j\in V}q(j).                \tag{5}
\]
Equations (4)--(5) rule out every strict inequality, and hence
\(f(j)=j\) for all \(j\).  Thus the unique observable support-compatible
matching is the true bijection \(\vartheta\).  It is a finite recovery map,
for example by enumerating the finitely many perfect matchings and retaining
the unique one satisfying \(\vartheta(c)\in S_c\).  Equation (3) now gives
the explicit graph recovery formula
\[
 \operatorname{Pa}(j)
 =S_{\vartheta^{-1}(j)}\setminus\{j\}.                  \tag{6}
\]

It remains to recover the numerical SEM parameters.  Put
\(P=\operatorname{Pa}(j)\), and let \(X\) be centered Gaussian with covariance
\(\Sigma_{c_0}\).  Define the proof intermediate \(\varepsilon=AX\).  The
baseline factorization gives, equation by equation,
\[
 \operatorname{Cov}(\varepsilon)
 =A\Sigma_{c_0}A^{\top}
 =A(A^{-1}\Omega A^{-\top})A^{\top}
 =\Omega.                                               \tag{7}
\]
Acyclicity makes \(B\) nilpotent and therefore
\[
 A^{-1}=(I_p-B)^{-1}=\sum_{m=0}^{p-1}B^m.               \tag{8}
\]
If \(i\in P\), then \(i\to j\).  For \(m\geq1\), a nonzero summand
contributing to \((B^m)_{ij}\) represents a directed path from \(j\) to
\(i\), which together with \(i\to j\) would be a directed cycle.  Also
\((I_p)_{ij}=0\).  Hence (8) implies
\((A^{-1})_{ij}=0\), and diagonality of \(\Omega\) gives
\[
 \operatorname{Cov}(X_i,\varepsilon_j)
 =e_i^{\top}A^{-1}\Omega e_j
 =\omega_j(A^{-1})_{ij}=0,
 \qquad i\in P.                                        \tag{9}
\]
The \(j\)-th row of \(\varepsilon=AX\) is
\[
 X_j=B_{jP}X_P+\varepsilon_j.                           \tag{10}
\]
Taking covariance of (10) with \(X_P\) and using (9) yields the normal
equations
\[
 (\Sigma_{c_0})_{Pj}
 =(\Sigma_{c_0})_{PP}B_{jP}^{\top}.                     \tag{11}
\]
Positive definiteness of \(\Sigma_{c_0}\) makes its principal submatrix
\((\Sigma_{c_0})_{PP}\) positive definite and therefore invertible.  Solving
(11) proves
\[
 B_{jP}^{\top}
 =(\Sigma_{c_0})_{PP}^{-1}(\Sigma_{c_0})_{Pj}.           \tag{12}
\]
Using (7), (9), and (10), or equivalently expanding the variance of the
residual in (10) and substituting (11)--(12), gives
\[
 \begin{aligned}
 \omega_j
 &=\operatorname{Var}(\varepsilon_j)\\
 &=(\Sigma_{c_0})_{jj}
   -(\Sigma_{c_0})_{jP}
    (\Sigma_{c_0})_{PP}^{-1}
    (\Sigma_{c_0})_{Pj}>0.                              \tag{13}
 \end{aligned}
\]
The strict inequality is also the positive Schur-complement property of the
positive-definite matrix \(\Sigma_{c_0}\).  If \(P=\varnothing\), the inverse
and products in (12)--(13) are interpreted as empty; then the coefficient row
has no parent entries and \(\omega_j=(\Sigma_{c_0})_{jj}>0\).

Finally consider the component matched to target \(j\).  If \(X^{(j)}\) has
its covariance and \(\eta^{(j)}=A^{(j)}X^{(j)}\), the perfect-surgery
factorization gives
\[
 \operatorname{Cov}(\eta^{(j)})
 =A^{(j)}\Sigma_j(A^{(j)})^{\top}
 =\Omega^{(j)}.                                         \tag{14}
\]
The \(j\)-th row of \(A^{(j)}\) is \(e_j^{\top}\), so
\(\eta_j^{(j)}=X_j^{(j)}\).  Taking the \((j,j)\) entry of (14) therefore
proves
\[
 \tau_j=(\Sigma_{\vartheta^{-1}(j)})_{jj}.              \tag{15}
\]

Equations (6), (12), (13), and (15) recover respectively \(G\), every
nonzero coefficient of \(B\), \(\omega\), and \(\tau\); zero entries of
\(B\) are fixed by the recovered parent sets.  The first step already
recovers the weight attached to each distinct covariance, so it also recovers
\(\pi\).  Thus the displayed operations define an observable recovery map on
the unordered covariance--weight collection, and the preceding equalities
prove, rather than define, that its value at every member of \(\mathcal M_p\)
is the generating
\(\theta_{\mathrm{sem}}=(c_0,\vartheta,G,B,\omega,\tau,\pi)\).  Composing
this map with the mixture identification in
\(\texttt{lem:component-distinctness}\) proves that \(P_{\mathrm{mix}}\)
uniquely determines \(\theta_{\mathrm{sem}}\), as claimed.
